\documentclass[notes,11pt, aspectratio=169, xcolor=table]{beamer}

\usepackage{pgfpages}
% These slides also contain speaker notes. You can print just the slides,
% just the notes, or both, depending on the setting below. Comment out the want
% you want.
\setbeameroption{hide notes} % Only slide
%\setbeameroption{show only notes} % Only notes
%\setbeameroption{show notes on second screen=right} % Both

\usepackage{helvet}
\usepackage[default]{lato}
\usepackage{array}

\newtheorem{proposition}{Proposition}

\usepackage{tikz}
\usetikzlibrary{shapes.geometric}
\usepackage{pgfplots}
\usepackage{graphicx}
\usepackage{verbatim}
\setbeamertemplate{note page}{\pagecolor{yellow!5}\insertnote}
\usetikzlibrary{positioning}
\usetikzlibrary{snakes}
\usetikzlibrary{calc}
\usetikzlibrary{arrows}
\usetikzlibrary{decorations.markings}
\usetikzlibrary{shapes.misc}
\usetikzlibrary{matrix,shapes,arrows,fit,tikzmark}
\usepackage{amsmath}
\usepackage{mathpazo}
\usepackage{hyperref}
\usepackage{lipsum}
\usepackage{multimedia}
\usepackage{graphicx}
\usepackage{multirow}
\usepackage{graphicx}
\usepackage{dcolumn}
\usepackage{bbm}
\usepackage[style=authoryear,sorting=nyt,uniquename=false]{biblatex}

\addbibresource{references.bib} 

\newcolumntype{d}[0]{D{.}{.}{5}}

\def\@@mybluebox[#1][#2]#3{
    \sbox\mytempbox{#3}%
    \mytemplen\ht\mytempbox
    \advance\mytemplen #1\relax
    \ht\mytempbox\mytemplen
    \mytemplen\dp\mytempbox
    \advance\mytemplen #2\relax
    \dp\mytempbox\mytemplen
    \colorbox{myblue}{\hspace{1em}\usebox{\mytempbox}\hspace{1em}}}


\usepackage{changepage}
\usepackage{appendixnumberbeamer}
\newcommand{\beginbackup}{
   \newcounter{framenumbervorappendix}
   \setcounter{framenumbervorappendix}{\value{framenumber}}
   \setbeamertemplate{footline}
   {
     \leavevmode%
     \hline
     box{%
       \begin{beamercolorbox}[wd=\paperwidth,ht=2.25ex,dp=1ex,right]{footlinecolor}%
%         \insertframenumber  \hspace*{2ex} 
       \end{beamercolorbox}}%
     \vskip0pt%
   }
 }
\newcommand{\backupend}{
   \addtocounter{framenumbervorappendix}{-\value{framenumber}}
   \addtocounter{framenumber}{\value{framenumbervorappendix}} 
}


\usepackage{graphicx}
\usepackage[space]{grffile}
\usepackage{booktabs}

% These are my colors -- there are many like them, but these ones are mine.
\definecolor{blue}{RGB}{0,114,178}
\definecolor{red}{RGB}{213,94,0}
\definecolor{yellow}{RGB}{240,228,66}
\definecolor{green}{RGB}{0,158,115}

\hypersetup{
  colorlinks=false,
  linkbordercolor = {white},
  linkcolor = {blue}
}


%% I use a beige off white for my background
\definecolor{MyBackground}{RGB}{255,253,218}

%% Uncomment this if you want to change the background color to something else
%\setbeamercolor{background canvas}{bg=MyBackground}

%% Change the bg color to adjust your transition slide background color!
\newenvironment{transitionframe}{
  \setbeamercolor{background canvas}{bg=yellow}
  \begin{frame}}{
    \end{frame}
}

\setbeamercolor{frametitle}{fg=blue}
\setbeamercolor{title}{fg=blue}
\setbeamertemplate{footline}[frame number]
\setbeamertemplate{navigation symbols}{} 
\setbeamertemplate{itemize items}{-}
\setbeamercolor{itemize item}{fg=blue}
\setbeamercolor{itemize subitem}{fg=blue}
\setbeamercolor{enumerate item}{fg=blue}
\setbeamercolor{enumerate subitem}{fg=blue}
\setbeamercolor{button}{bg=MyBackground,fg=blue,}



% If you like road maps, rather than having clutter at the top, have a roadmap show up at the end of each section 
% (and after your introduction)
% Uncomment this is if you want the roadmap!
% \AtBeginSection[]
% {
%    \begin{frame}
%        \frametitle{Roadmap of Talk}
%        \tableofcontents[currentsection]
%    \end{frame}
% }
\setbeamercolor{section in toc}{fg=blue}
\setbeamercolor{subsection in toc}{fg=red}
\setbeamersize{text margin left=1em,text margin right=1em} 

\newenvironment{wideitemize}{\itemize\addtolength{\itemsep}{10pt}}{\enditemize}

\usepackage{environ}
\NewEnviron{videoframe}[1]{
  \begin{frame}
    \vspace{-8pt}
    \begin{columns}[onlytextwidth, T] % align columns
      \begin{column}{.58\textwidth}
        \begin{minipage}[t][\textheight][t]
          {\dimexpr\textwidth}
          \vspace{8pt}
          \hspace{4pt} {\Large \sc \textcolor{blue}{#1}}
          \vspace{8pt}
          
          \BODY
        \end{minipage}
      \end{column}%
      \hfill%
      \begin{column}{.42\textwidth}
        \colorbox{green!20}{\begin{minipage}[t][1.2\textheight][t]
            {\dimexpr\textwidth}
            Face goes here
          \end{minipage}}
      \end{column}%
    \end{columns}
  \end{frame}
}

\title[]{International Trade: Lecture 1}
\subtitle[]{Preliminaries}
\author[Góes]
{Carlos Góes\inst{1}}
\date{Fall 2025}
\institute[GWU]{\inst{1} George Washington University }



\begin{document}

%%% TIKZ STUFF
\tikzset{   
        every picture/.style={remember picture,baseline},
        every node/.style={anchor=base,align=center,outer sep=1.5pt},
        every path/.style={thick},
        }
\newcommand\marktopleft[1]{%
    \tikz[overlay,remember picture] 
        \node (marker-#1-a) at (-.3em,.3em) {};%
}
\newcommand\markbottomright[2]{%
    \tikz[overlay,remember picture] 
        \node (marker-#1-b) at (0em,0em) {};%
}
\tikzstyle{every picture}+=[remember picture] 
\tikzstyle{mybox} =[draw=black, very thick, rectangle, inner sep=10pt, inner ysep=20pt]
\tikzstyle{fancytitle} =[draw=black,fill=red, text=white]
%%%% END TIKZ STUFF



%----------------------------------------------------------------------%
%-------------------       TITLE PAGE       ---------------------------%
%----------------------------------------------------------------------%





%----------------------------------------------------------------------%






%----------------------------------------------------------------------%
%----------------------------------------------------------------------%

%----------------------------------------------------------------------%
\frame{\titlepage}
\addtocounter{framenumber}{-1}
%----------------------------------------------------------------------%



%----------------------------------------------------------------------%
%----------------------------------------------------------------------%

\section{Housekeeping}

\begin{frame}{Outline}
\begin{wideitemize}
\item About me
\item Course objectives and organization
\item Why do we trade: the Principle of Comparative Advantage
    
\end{wideitemize}

\end{frame}


\begin{frame}{About me}

\begin{columns}[T] % align columns
\begin{column}{.38\textwidth}
I'm originally from Brazil...
\vspace{10cm}
  \makebox[\linewidth][c]{
    \resizebox{\linewidth}{!}{
    \includegraphics{figs/brazil.png}
      }
    }
\end{column}%
\hfill%
\begin{column}{.58\textwidth}
  \begin{wideitemize}
    \item Education:
    \begin{itemize}
        \item Ph.D. \& M.A., Economics - UC San Diego 
        \item M.A.: Int'l Econ + Int'l Relations, Johns Hopkins 
    \end{itemize}
    \item Professional:
    \begin{itemize}
        \item Economist, World Bank Group 
        \item Research Economist, World Trade Organization
        \item Senior Economic Advisor, President of Brazil
    \end{itemize}    
    \item Misc:
    \begin{itemize}
        \item Twitter: @goescarlos
        \item Website: www.carlosgoes.com
    \end{itemize}    
    \item Research: Macro \& International Trade
  \end{wideitemize}
\end{column}%
\end{columns}
\end{frame}


\begin{frame}{Course objectives}
\begin{wideitemize}
\item International Trade (xxx).
\item Course objectives:
    \begin{enumerate}
        \item Introduction to modern trade theory + trade data
        \item Understand how theory + data inform policy
        \item Critically assess the policy debates
    \end{enumerate}

\item What do you need to know?
    \begin{enumerate}
        \item Pre-reqs: \textbf{xxxx}
        \item Most of you are econ majors, so I expect you to have heard of:
        \begin{enumerate}
            \item Indifference curves
            \item Production functions / production possibility sets
            \item Equilibrium
        \end{enumerate}
        \item I expect you know how to take a derivative (or be willing to re-learn how to do it)
        \item This is \textbf{not} a mathematical economics course: no long, detailed, rigorous proofs
    \end{enumerate}

\end{wideitemize}    
\end{frame}


\begin{frame}{Textbook}

\begin{columns}[T] % align columns
\begin{column}{.4\textwidth}
  \makebox[\linewidth][c]{
    \resizebox{\linewidth}{!}{
    \includegraphics{figs/KOM-IT.png}
      }
    }
\end{column}%
\hfill%
\begin{column}{.58\textwidth}
  \begin{wideitemize}
  \item \textit{International Trade: Theory and Policy.}
    \item Edition: 11th
    \item Authors:
  \begin{itemize}
    \item Paul Krugman (Nobel)
    \item Maurice Obstfeld (former IMF Chief)
    \item Marc Melitz (future Nobel)
  \end{itemize}
  \end{wideitemize}
\end{column}%
\end{columns}
\end{frame}


\begin{frame}{Grading}
\begin{wideitemize}
    \item 
\end{wideitemize}    
\end{frame}

\begin{frame}{Key dates}
\begin{wideitemize}
    \item 
\end{wideitemize}    
\end{frame}

\begin{frame}{Policies}
\begin{wideitemize}
    \item \textcolor{blue}{Attendance}
    \begin{itemize}
        \item I do not require attendance
        \item You are an adult and I will treat you as such
        \item If you skip all the classes and do well in my class, congrats (low probability)
        \item If you attend all classes and follow the slides, you will get a good grade (high probability) 
    \end{itemize}

\onslide<2->{
    \item \textcolor{blue}{Technology}
    \begin{itemize}
        \item For take-home exercises, you can use whatever tools you want, including AI
        \item However, you need to type/write the answers yourself
        \item In other words, you can't copy and paste images of ChatGPT into the answers
        \item AI can be your friend (I encourage it) but it's an assistant not a replacement for effort
    \end{itemize}
}
\onslide<3->{
     \item \textcolor{blue}{Honesty and regrading policy}
    \begin{itemize}
        \item The main rule is: do not be dishonest (you will get caught)
        \item If you have a problem, let me know in advance -- I am reasonable
        \item Regrade policy: arithmetical error (I summed scores wrong); or full exam regrade (your score might go up or down)
        \item You cannot appeal on a single question
        \end{itemize}   
}
\end{wideitemize}    
\end{frame}


\begin{frame}{Data + coding}
\begin{wideitemize}
    \item \textcolor{blue}{Coding}
    \begin{itemize}
        \item Aside from trade knowledge, I want you to get skills out of this class
        \item I will ask you to download data, transform it using computer code, plot charts and tables
        \item If you don't know how, don't freak out... I'll help you (also, use Google + AI). Learn!
    \end{itemize}
\onslide<2->{
    \item \textcolor{blue}{Python}
    \begin{itemize}
        \item Python is a general purpose language... people actually code apps with it
        \item It is also very useful for data science:
        \begin{itemize}
            \item downloading data
            \item transforming data
            \item using statistics to summarize or correlate it
            \item plotting charts and maps
        \end{itemize}
    \end{itemize}
}
\onslide<3->{
    \item \textcolor{blue}{Anaconda}
    \begin{itemize}
        \item Free + open source distribution
        \item Optimized for data science
        \item \url{https://www.anaconda.com/download}
        \item Installing + exploring lab next class
    \end{itemize}
}
\end{wideitemize}    
\end{frame}

\end{document}
